\documentclass{article}
\usepackage{amsmath,amssymb}
% Copyright (c) 2026 Alberto Villa Osorno.
% SPDX-License-Identifier: MIT
% Shared notation only; this file is included by buildable math documents.

\newcommand{\TritSet}{\{0,1,2\}}
\newcommand{\WordModulus}[1]{3^{#1}}
\newcommand{\WordDomain}[1]{\mathcal{W}_{#1}}
\newcommand{\MemoryState}[1]{\mathcal{M}_{#1}}
\newcommand{\CrazyDigit}{\gamma}
\newcommand{\CrazyN}[3]{\Gamma_{#1}\!\left(#2,#3\right)}
\newcommand{\RotateN}[2]{\rho_{#1}\!\left(#2\right)}
\newcommand{\DecodeOp}{\delta}
\newcommand{\EncryptOp}{\varepsilon}
\newcommand{\LoadOp}{\mathsf{Load}}
\newcommand{\LowerOp}{\mathsf{Lower}}
\newcommand{\VerifyOp}{\mathsf{Verify}}
\newcommand{\ObservedEquivalent}[1]{\equiv_{\mathrm{obs},#1}}
\newcommand{\CostVector}[1]{\mathbf{K}\!\left(#1\right)}
\newcommand{\Transition}[1]{\xrightarrow{#1}}


\title{Versioned Malbolge Profile Mathematics}
\author{Malbolge Project}
\date{2026-07-26}

\begin{document}
\maketitle

\section{Scope}

This document fixes mathematical notation shared by versioned Malbolge target
profiles. A profile chooses a positive ternary word width $N$ and therefore the
word modulus
\[
  W_N = \WordModulus{N}.
\]
The profile-specific implementation contract supplies finite translation tables
and the remaining non-mathematical source-admission details.

\section{Word and memory domains}

Define
\[
  \WordDomain{N} = \{0,1,\ldots,W_N-1\}.
\]
Every word $x\in\WordDomain{N}$ has the unique expansion
\[
  x = \sum_{i=0}^{N-1} x_i3^i,
  \qquad x_i\in\TritSet.
\]
For the single-word-modular memory model, memory is a total function
\[
  M : \WordDomain{N}\to\WordDomain{N},
\]
and registers satisfy
\[
  A,C,D\in\WordDomain{N}.
\]
Pointer succession is
\[
  \operatorname{succ}_N(x)=(x+1)\bmod W_N.
\]

\section{Rotation}

The profile-width ternary right rotation is
\[
  \RotateN{N}{x}
    = \left\lfloor\frac{x}{3}\right\rfloor
      + (x\bmod 3)3^{N-1}.
\]
It maps $\WordDomain{N}$ to itself and moves the least-significant trit to the
most-significant position.

\section{Crazy operation}

Define the digit operation
\[
  \CrazyDigit : \TritSet\times\TritSet\to\TritSet
\]
by
\[
\begin{array}{c|ccc}
\CrazyDigit(d,a) & a=0 & a=1 & a=2 \\
\hline
 d=0 & 1 & 0 & 0 \\
 d=1 & 1 & 0 & 2 \\
 d=2 & 2 & 2 & 1
\end{array}
\]
with data trit $d$ first and accumulator trit $a$ second. The word operation is
\[
  \CrazyN{N}{d}{a}
    = \sum_{i=0}^{N-1}\CrazyDigit(d_i,a_i)3^i.
\]
This orientation is repository-wide notation; optimized tables may not silently
transpose the two operands.

\section{Decode and self-modification}

Let $X_1$ and $X_2$ denote the immutable 94-entry decode and encryption tables.
For graphical cells $m\in\{33,\ldots,126\}$,
\[
  \DecodeOp(m,c)
    = X_1\!\left[(m-33+c)\bmod 94\right],
\]
and
\[
  \EncryptOp(m)=X_2[m-33].
\]
A non-halting transition first applies its instruction-specific effect, then
requires the resulting encryption target $M[C]$ to remain graphical, applies
$M[C]\leftarrow\EncryptOp(M[C])$, and only then advances code and data pointers
with $\operatorname{succ}_N$. A jump that changes $C$ therefore also changes the
encryption target. Rejected encryption is atomic.

\section{I/O and termination}

Byte input maps an available byte $b\in\{0,\ldots,255\}$ to $A\leftarrow b$.
End of input maps to the all-two-trit sentinel
\[
  A\leftarrow W_N-1.
\]
Byte output emits $A\bmod256$. The halt instruction terminates before encryption
or pointer advancement. Fetching a non-graphical current code cell terminates
immediately under the normative profile semantics.

\section{Loading notation}

For a profile $p$ with width $N$, write
\[
  \LoadOp_p(s)=M_0
\]
when source bytes $s$ satisfy the profile source-admission rules and determine an
exact initial memory image $M_0\in\MemoryState{N}$. Remaining cells are generated
by the profile-width crazy recurrence from the previous two loaded/generated
words. Loader rejection is outside the execution transition relation.

\section{Execution and observable equivalence}

Let one complete guest state be
\[
  S=(A,C,D,M,I,O),
\]
where $I$ is the remaining byte-input stream and $O$ is committed byte output.
Write
\[
  S \Transition{p} S'
\]
for one committed normative step under profile $p$.

For programs $P$ and $Q$, write
\[
  P\ObservedEquivalent{p}Q
\]
when, for the declared input domain, both executions have identical observable
byte I/O, termination, and all verifier-required state observations under the
same profile. An optimization or native tier may use another representation but
must discharge this equivalence obligation over its admitted domain.

\section{Compiler and research notation}

Compiler lowering under profile $p$ is written
\[
  C \Transition{\LowerOp_p} P,
\]
while independent acceptance of candidate $P$ is written
\[
  \VerifyOp(P,p)=1.
\]
Research cost vectors use
\[
  \CostVector{P}
    = (K_{\mathrm{size}},K_{\mathrm{time}},K_{\mathrm{memory}},\ldots).
\]
Cost ordering never implies semantic correctness; candidates enter comparative
optimization only after the declared verifier/equivalence obligation succeeds.

\section{Canonical specializations}

The written historical profile uses $N=10$, giving $W_{10}=59049$. The current
`malbolge-2026.2` profile uses $N=14$, giving
\[
  W_{14}=4782969.
\]
These are specializations of one mathematical profile model, not permission to
reinterpret historical resource limits as current-language constants.

\end{document}
